\documentclass[a4paper,fleqn]{cas-dc}

\usepackage[authoryear,longnamesfirst]{natbib}

% --- Additional packages used by the article body (not already
%     provided by cas-dc.cls / cas-common.sty) --------------------
\usepackage{enumitem}   % \begin{itemize}[leftmargin=*] in future_work.tex
\usepackage{bbm}        % \mathbbm{1} indicator function in tto.tex

% --- Custom macros (carried over from the original preamble.tex) ---
\newcommand{\R}{\mathbb{R}}
\DeclareMathOperator{\diag}{diag}
\DeclareMathOperator*{\argmin}{arg\,min}
\newcommand{\SAE}{\mathrm{SAE}}
\newcommand{\DAE}{\mathrm{DAE}}

\begin{document}
\let\WriteBookmarks\relax
\def\floatpagepagefraction{1}
\def\textpagefraction{.001}

% Short title / authors for running heads
\shorttitle{Refinements to the SAE--TS Framework for Email Outreach}
\shortauthors{T. Ivanov}

% Main title
\title[mode=title]{Refinements to the SAE--Thompson Sampling Framework for
  Active Email Outreach}

% First (and only) author
\author[1]{Tymofii Ivanov}
\affiliation[1]{organization={Faculty of Information Technology, Czech
    Technical University in Prague},
            city={Prague},
            country={Czech Republic}}

% Unnumbered note reproducing the "Extends:" note from the original draft
\nonumnote{\emph{Extends:} \v{Z}id, \v{C}., Alves, R., and Kord\'ik, P.\ 2025.
  \emph{Active Recommendation for Email Outreach Dynamics.} In
  \emph{Proceedings of the 34th ACM International Conference on Information
  and Knowledge Management (CIKM~'25)}, Seoul, Republic of Korea.\
  \url{https://doi.org/10.1145/3746252.3760832}}

\begin{abstract}
\v{Z}id, Alves, and Kord\'ik (2025) frame cold-start recipient selection for a new email
campaign as an iterative multi-armed bandit driven by Thompson Sampling
(TS), with parameters derived from a shallow autoencoder (SAE) trained on
historical recipient--template opens. Several components of that model---a
fixed exploration/exploitation coefficient, an unweighted historic open
rate, a linear aggregation for the current-state confidence term, training
on binary labels only, and a shallow autoencoder architecture---were fixed
by the original authors for simplicity. We revisit each of these choices
and propose concrete, mathematically grounded refinements: dynamic
scheduling of the blend coefficient, a recency-weighted historic rate, a
nonlinear forward-pass confidence term, a variance-based influence score, a
factored score function, two ways of exposing the autoencoder to
previously unused time-to-open (TTO) signal, and a deep, nonlinear
autoencoder. We report experimental results for each proposal against a
reproduced baseline and outline directions for further work.
\end{abstract}

\begin{highlights}
\item Revisits five simplifying design choices of the SAE--Thompson
  Sampling email-outreach bandit of \citet{zid2025active}.
\item A variance-based influence score gives the largest low-sending-
  fraction recall gain (0.479 vs.\ 0.385 baseline at 5\%).
\item Recency-weighted historic rate and a factored score also raise
  Recall-AUC (0.903 and 0.900 vs.\ 0.894 baseline).
\item A deep, nonlinear autoencoder improves recall modestly but with
  markedly higher variance, indicating a need for regularization.
\end{highlights}

\begin{keywords}
Thompson sampling \sep multi-armed bandits \sep active learning \sep
email recommendation \sep autoencoders \sep collaborative filtering
\end{keywords}

\maketitle

\section{Introduction}
\label{sec:introduction}

\citet{zid2025active} frame cold-start recipient selection for a
new email campaign as an iterative multi-armed bandit: each arm is a
recipient, and Thompson Sampling (TS) parameters are derived from a shallow
autoencoder (SAE) trained on historical recipient--template opens, avoiding
retraining during the active-learning phase. Several components of that
model --- a fixed exploration/exploitation coefficient, an unweighted
historic open rate, and a linear aggregation for the current-state
confidence term --- were fixed by the original authors for simplicity. We
revisit each of these choices and propose concrete, mathematically grounded
refinements: dynamic scheduling of the blend coefficient, a
recency-weighted historic rate, a nonlinear forward-pass confidence term, a
variance-based influence score, and a factored score function. We also
describe the evaluation protocol used to compare each proposal against a
reproduced baseline.

\section{Recap of the Original Model}
\label{sec:recap}

\paragraph{Setup.}
We observe $X \in \{0,1\}^{n \times m}$, where $n$ is the number of
templates, $m$ the number of recipients, and $X_{i,j}=1$ if recipient $j$
opened template $i$. For a new template $n{+}1$, sent over $b$ batches
within a deadline $T$, the model selects at each batch a small subset of
recipients so as to maximize eventual openers reached early.

\paragraph{Shallow autoencoder.}
A SAE with encoder/decoder $E, D \in \R^{m\times d}$ is trained as
\begin{equation*}
\min_{E,D} \ell\big(X,\ \sigma(XB_{E,D})\big),
\qquad
B_{E,D} = ED^\top - \diag\big([E\odot D]\mathbf{1}_m\big),
\end{equation*}
with $\ell$ the element-wise binary cross-entropy and the diagonal
constraint ruling out trivial self-prediction. We write
$f_{\SAE}(x) = \sigma(x^\top B_{E,D})$ for the trained forward pass and
$\Sigma = \sigma(B_{E,D})$; entry $\Sigma_{i,j}$ is the probability of
recipient $j$ opening a template given that only recipient $i$ has opened
it --- a measure of $i$'s influence on $j$.

\paragraph{Arm score.}
The bandit score for recipient $j$ at time $t$ is
\begin{equation}
\label{eq:score}
s_j(t) = \alpha\,\phi_j p_j + (1-\alpha) f_j(t),
\end{equation}
\begin{equation}
\label{eq:phi-p-f}
\phi_j = \frac1n\sum_{i=1}^n X_{i,j},
\qquad
p_j = \frac{1}{m-1}\sum_{i\ne j}\Sigma_{j,i},
\qquad
f_j(t) = \bar x(t)^\top \Sigma_{:j},
\end{equation}
where $\phi_j$ is recipient $j$'s historic open rate, $p_j$ is a general
influence score, $f_j(t)$ is a confidence score specific to template
$n{+}1$ given the observed (normalized) state $\bar x(t)$, and
$\alpha \in [0,1]$ trades off the two. TS parameters are then
$\alpha_j(t) = Gs_j(t)$, $\beta_j(t) = G(1-s_j(t))$ for a confidence
modifier $G$.

Two orthogonal readings of Eq.~\eqref{eq:score} motivate our refinements:
$\phi_j p_j$ vs.\ $f_j(t)$ can be read as \emph{exploration vs.\
exploitation}, or equivalently as \emph{historic vs.\ current} data.

\section{Proposed Refinements}
\label{sec:refinements}

\subsection{Dynamic $\alpha$-Scheduling}
\label{sec:alpha-scheduling}

A fixed $\alpha$ ignores that the value of exploration should shrink as the
campaign progresses, echoing the broader practice of decaying exploration
coefficients over time in the bandit and reinforcement-learning literature
\citep{auer2002finite}. We schedule $\alpha(t)$ against two natural progress
variables: $\mu = N(t)/N$, the fraction of the sending budget already used,
and $\tilde o = o(t)/m$, the fraction of recipients who have already
opened, with $o(t) = x(t)^\top\mathbf{1}$.

\paragraph{Send-volume-based (function of $\mu$).}
A \textbf{linear} schedule from $l$ to $r$, in the spirit of the linearly
annealed $\varepsilon$-greedy exploration rate used to train DQN agents
\citep{mnih2015human},
\begin{equation}
\label{eq:alpha-linear}
\alpha(\mu) = l(1-\mu) + r\mu,
\end{equation}
and a \textbf{geometric} (log-linear) schedule, analogous to the geometric
cooling schedules of simulated annealing \citep{mahdi2017performance},
\begin{equation}
\label{eq:alpha-geometric}
\alpha(\mu) = l^{1-\mu}\, r^{\mu}.
\end{equation}

\paragraph{Open-rate-based (function of $\tilde o$).}
Since $f_j(t)$ is itself a point estimate that becomes more reliable as
more opens are observed, we also consider a \textbf{confidence} schedule,
hyperbolic in $\tilde o$ --- paralleling how the confidence radius of
UCB-style bandit policies shrinks as observations accumulate
\citep{auer2002finite}:
\begin{equation}
\label{eq:alpha-confidence}
\alpha(\tilde o) = \frac{\kappa}{\tilde o + \kappa},
\end{equation}
with dampening factor $\kappa$ (so $\alpha(\kappa)=\tfrac12$). All three
reduce to the constant baseline at their degenerate limits ($l=r$, or
$\kappa\to\infty$). A schedule jointly conditioned on both $\mu$ and
$\tilde o$ is conceptually natural but is left to future work
to avoid an unconstrained hyperparameter count.

\subsection{Recency-Weighted Historic Rate $\phi_j$}
\label{sec:recency-phi}

The original $\phi_j$ weights every past template equally, ignoring drift
in recipient behavior over time --- a limitation well studied in
time-aware collaborative filtering \citep{ding2005time,koren2009collaborative}.
Assuming templates
$X_{1:},\dots,X_{n:}$ were sent at uniform intervals, we assign template
$i$ an exponential weight, following the time-weighting scheme of
\citet{ding2005time},
\begin{equation}
\label{eq:omega}
\omega_i = 2^{-d_i/h}, \qquad d_i = \frac{n-i}{n},
\end{equation}
with half-life $h>0$, and redefine
\begin{equation}
\label{eq:phi-weighted}
\phi_j = \frac{\omega^\top X_{:j}}{\sum_i \omega_i}.
\end{equation}
Normalizing by $\sum_i\omega_i$ keeps $\phi_j\in[0,1]$. As $h\to\infty$
every $\omega_i\to1$, recovering the original unweighted average as a
special case.

\subsection{Forward-Pass Definition of $f_j(t)$}
\label{sec:forward-f}

$f_j(t) = \bar x(t)^\top\Sigma_{:j}$ is a linear combination of individual
influences $\Sigma_{i,j}$ and therefore captures only the \emph{average}
effect of openers, never their joint interaction. We instead pass the full
current state through the trained network, in the spirit of nonlinear
autoencoder-based collaborative-filtering predictors such as AutoRec
\citep{sedhain2015autorec}:
\begin{equation}
\label{eq:f-forward}
f_j(t) = \sigma\big(x(t)^\top B_{E,D}\big)e_j = f_{\SAE}(x(t))\,e_j.
\end{equation}
Because $\sigma$ is nonlinear, $\sigma(x(t)^\top B_{E,D}) \ne
x(t)^\top\sigma(B_{E,D})$ in general (equality holds only for one-hot
$x(t)$), so Eq.~\eqref{eq:f-forward} is a genuinely different,
jointly-conditioned probability estimate rather than a relabeling of
Eq.~\eqref{eq:phi-p-f}.

\subsection{Negative Feedback via a Second SAE}
\label{sec:negative-feedback}

The model of Section~\ref{sec:recap} learns from opens alone: both the
training matrix $X$ and the state $x(t)$ encode a recipient who was
contacted but did not open exactly as one who was never contacted. The data,
however, distinguish three cases --- \emph{sent and opened}, \emph{sent but
not opened}, and \emph{not sent} --- and the middle one is informative: an
ignored email is evidence against opening for similar recipients. We
recover it through the send matrix $S\in\{0,1\}^{n\times m}$, where
$S_{i,j}=1$ if template $i$ was sent to recipient $j$ (so $X_{i,j}\le
S_{i,j}$), and train a second SAE on the complementary signal.

\paragraph{Negative channel.}
We write $X^{+}=X$ and define the sent-but-unopened matrix
\begin{equation}
\label{eq:x-neg}
X^{-}_{i,j} = S_{i,j}\,\big(1-X_{i,j}\big)\in\{0,1\}.
\end{equation}
We train two SAEs, $(E^{+},D^{+})$ on $X^{+}$ and $(E^{-},D^{-})$ on
$X^{-}$, each with the objective of Section~\ref{sec:recap} (same bottleneck
size $d$), and write $B^{\pm}=B_{E^{\pm},D^{\pm}}$; the positive SAE is the
original one. At time $t$ the state of template $n{+}1$ splits accordingly:
$x^{+}(t)=x(t)$ marks recipients who have opened, and $x^{-}(t)\in\{0,1\}^m$
marks recipients who were sent the template in $[0,t)$ and have not opened
it. The three kinds of recipient are thus separated, with
$(x^{+}_j(t),x^{-}_j(t))$ equal to $(1,0)$, $(0,1)$ and $(0,0)$ for
opened, sent-but-unopened and not-yet-contacted recipients, respectively.

\paragraph{Logit blending.}
Let
\begin{equation}
\label{eq:z-pm}
z^{+}_j(t) = x^{+}(t)^\top B^{+}e_j,
\qquad
z^{-}_j(t) = -\,x^{-}(t)^\top B^{-}e_j
\end{equation}
be the positive and negative logits. The minus sign encodes the direction
of the evidence: a large $B^{-}_{i,j}$ means that $i$ failing to open makes
$j$ failing more likely, which must \emph{lower} $f_j(t)$. Equivalently, the
negative SAE is queried with the state $-x^{-}(t)$ (a $-1$ for every
sent-but-unopened recipient), while it is trained on the $0/1$ targets
$X^{-}$ so that the cross-entropy remains well-defined. We blend the two
channels in logit space,
\begin{equation}
\label{eq:f-neg}
f_j(t) = \sigma\big(\lambda z^{+}_j(t) + (1-\lambda)z^{-}_j(t)\big),
\qquad \lambda\in[0,1].
\end{equation}
Since the blend acts before $\sigma$, Eq.~\eqref{eq:f-neg} extends the
forward-pass definition of Section~\ref{sec:forward-f}, not the linear form
of Eq.~\eqref{eq:phi-p-f}, which has no logit to blend. For $\lambda=1$ it
reduces exactly to Eq.~\eqref{eq:f-forward}; with no evidence
($x^{+}(t)=x^{-}(t)=\mathbf{0}$) it returns the neutral prior $\tfrac12$ of
the original SAE; and $\lambda=\tfrac12$ gives the unweighted sum of logits
up to a global rescaling by $2$.

\paragraph{Role of $\lambda$.}
A weight is needed because the two channels are unbalanced. With open rates
around $9\%$ in the original dataset, sent-but-unopened entries outnumber
opens roughly ten to one, so $X^{-}$ is far denser than $X^{+}$ and
$x^{-}(t)$ typically has far more non-zeros than $x^{+}(t)$; an unweighted
$z^{-}_j(t)$ would accumulate many more terms than $z^{+}_j(t)$ and
dominate. Moreover, a missing open is weaker evidence than an open: the
recipient may simply not have read the email yet. The coefficient $\lambda$
trades the two off, and, like $\alpha$, is a hyperparameter of the model.
Both SAEs are trained once offline, so the active-learning phase remains
free of retraining, and $f_j(t)$ from Eq.~\eqref{eq:f-neg} can replace
$f_j(t)$ in Eq.~\eqref{eq:score} or in $\pi_j(t)$ of
Eq.~\eqref{eq:pi-u} without further changes.

\subsection{Variance-Based Definition of $p_j$}
\label{sec:variance-p}

$p_j$ is meant to reward recipients whose opening is \emph{informative}
about others, in the sense of informativeness used throughout active
learning. The original definition,
$p_j = \frac{1}{m-1}\sum_{i\ne j}\Sigma_{j,i}$, rewards recipients that
raise the \emph{average level} of other recipients' $f$-scores --- which is
not the same as reducing uncertainty about them. We instead reward maximal
\emph{spread} induced across other recipients' scores, following
variance-based exploration bonuses in the multi-armed bandit literature
\citep{audibert2009exploration}:
\begin{equation}
\label{eq:p-variance}
p_j = \frac{4}{m-1}\sum_{i\ne j}\big(\Sigma_{j,i} - \bar\Sigma_{j:}\big)^2,
\qquad
\bar\Sigma_{j:} = \frac{1}{m-1}\sum_{i\ne j}\Sigma_{j,i}.
\end{equation}
This is the variance of $\Sigma_{j,:}$ (excluding $j$), scaled by $4$ so
that the maximum possible variance of a Bernoulli variable (0.25) maps to
1, keeping $p_j \in [0,1]$.

\subsection{Factored Score $s_j(t) = \pi_j(t)\,u_j(t)$}
\label{sec:factored-score}

The linear combination in Eq.~\eqref{eq:score} conflates the two
orthogonal axes noted in Section~\ref{sec:recap} --- it is not clear, for
instance, why exploration should draw solely on historic data while
exploitation draws solely on current data. We instead factor the score
into a \emph{probability-of-opening} term and a \emph{utility-of-opening}
term, each independently schedulable:
\begin{equation}
\label{eq:pi-u}
\pi_j(t) = \alpha(t)\phi_j + (1-\alpha(t))f_j(t),
\qquad
u_j(t) = \beta(t)p_j + (1-\beta(t))\cdot 1,
\end{equation}
\begin{equation}
\label{eq:s-factored}
s_j(t) = \pi_j(t)\,u_j(t)
       = \Big(\alpha(t)\phi_j + [1-\alpha(t)]f_j(t)\Big)
         \Big(\beta(t)p_j + [1-\beta(t)]\Big).
\end{equation}
Here $\pi_j(t)$ blends historic and current opening-probability estimates,
while $u_j(t)$ blends the direct value of an open (the constant $1$) with
its indirect information value $p_j$. This factorization reuses every
primitive introduced above unchanged, so refinements from
Sections~\ref{sec:alpha-scheduling}--\ref{sec:variance-p} can be dropped
into $\pi_j(t)$ and $u_j(t)$ directly --- we use a confidence schedule
(Eq.~\eqref{eq:alpha-confidence}) for $\alpha(t)$ and a send-volume
schedule (Eq.~\eqref{eq:alpha-linear}/\eqref{eq:alpha-geometric}) for
$\beta(t)$.


\section{Experiments}
\label{sec:experiments}

\subsection{Methodology}
\label{sec:methodology}

For each variant we grid-search its hyperparameters and select the
configuration maximizing validation-set Recall-AUC,
$\mathrm{AUC} = \int_0^1 \mathrm{Recall}(\tau)\,d\tau$, where
$\mathrm{Recall}(\tau)$ is recall at sending fraction $\tau$. Variants are
compared on test-set Recall@5/15/25/35\% and Recall-AUC, as mean $\pm$ std
over repeated simulations. Before testing modifications we reproduced the
original model on our pipeline; at the operating points reported by
\citet{zid2025active} (25/50/75\%) our reproduction matches the published
point estimates exactly ($0.923$, $0.975$, $0.989$), with larger variance
attributable to fewer simulation runs.

\section{Future Work}
\label{sec:future-work}

\begin{itemize}[leftmargin=*]

\item \textbf{Combined $\alpha$-scheduling.} A schedule jointly
  conditioned on send-progress $\mu$ and open-progress $\tilde o$ rather
  than either alone (Section~\ref{sec:alpha-scheduling}); the main
  obstacle is the added hyperparameter count.

\item \textbf{Plug-and-play composition under the factored score.}
  Systematically combine the primitive-level refinements ---
  recency-weighted $\phi_j$, forward-pass $f_j(t)$, variance-based $p_j$,
  TTO- or deep-autoencoder-derived $\Sigma$ --- inside $\pi_j(t)$ and
  $u_j(t)$ (Eq.~\eqref{eq:pi-u}) rather than evaluating each in isolation,
  and search jointly rather than ablating one factor at a time, since the
  gains in Section~\ref{sec:discussion} are not guaranteed to be additive.

\item \textbf{Hyperbolic decay for $Y$.}
  $Y_{i,j} = \kappa_\delta/(\Delta_{i,j}+\kappa_\delta)$ gives a heavier
  tail than the exponential label of Eq.~\eqref{eq:Y-decay} and merits a
  direct comparison.

\item \textbf{TTO-to-TTO autoencoder.} Replacing the SAE input with $Y$ as
  well, $\min_{E,D}\ell(Y,\sigma(YB_{E,D}))$, exposes TTO on both sides but
  changes the role of $e_j$ --- an instantaneous open becomes a boundary
  point rather than a typical input --- so $p_j$ can no longer be queried
  at $e_j$ and would need re-deriving, e.g.\ via a recipient- or
  population-level mean-TTO query, with unclear risk of double-counting
  against $\phi_j$.

\item \textbf{Asymmetric deep encoder/decoder.} The two-layer network of
  Section~\ref{sec:deep-ae} uses matched encoder/decoder depth and width;
  decoupling them (e.g.\ a wider encoder with a narrower, more
  constrained decoder, or vice versa) may better balance representational
  capacity against the denoising-based regularization, and warrants
  exploration alongside a wider sweep over masking rate, dropout, and
  bottleneck width $d$ to address the variance observed in
  Section~\ref{sec:discussion}.

\end{itemize}


\bibliographystyle{cas-model2-names}
\bibliography{references}

\end{document}
